% !TeX root = ../教材题库_全部条目_转写工作稿.tex
% Source: XJ-B1-C05-SU-P205-Y1165
% 第5章 / 小结与复习 小结与复习 / 教材 p.205 / PDF 第212页
% OCR未校对：公式与文字均可能有误；未完成数学TeX转写；禁止直接发给学生
% 原题及相邻原解截图（核对用）：../题目原图/XJ-B1-C05-SU-P205-Y1165-01.webp
% 本文件只在转写完成后作为题干使用；OCR 条目尚不是可靠题干。
\begin{bankitem}{第5章 小结与复习 第4题}{第5章 / 小结与复习 小结与复习 / 教材 p.205 / PDF 第212页}{XJ-B1-C05-SU-P205-Y1165}{OCR未校对：公式与文字均可能有误；未完成数学TeX转写；禁止直接发给学生}
\begin{ocrdraft}
[教材p.205 / PDF第212页]
\par
4．三角函数的图象与性质是本章的重点．试借助图象归纳总结正弦函数、余
\par
弦函数和正切函数的性质．与我们前面所学的函数知识相比，这些三角函数有哪
\par
些不同的性质？
\end{ocrdraft}
\end{bankitem}
